% ═══════════════════════════════════════════════════════════════
%  chapters/ch02-flashing.tex
%  Прошивка STM32 из командной строки
% ═══════════════════════════════════════════════════════════════

\chapter{Прошивка STM32 из командной строки}
\label{ch:flashing}

\epigraph{<<Первая программа~--- всегда самая сложная.
Не потому, что она сложная, а потому, что нужно
настроить \emph{всё} окружение~--- и ни разу не
ошибиться>>.}{}

\section{Введение}
\label{sec:flashing-intro}

В \texttt{ch01} мы \emph{убедились}, что у нас есть
\emph{всё} необходимое: Nucleo, Blue Pill, ST-Link,
компьютер на Linux. Но пока мы \emph{ничего} не
прошили. Платы лежат на столе, светодиоды мигают
\emph{заводской} прошивкой, а мы~--- не понимаем,
\emph{что} там внутри.

Пришло время \textbf{прошить} первую программу.

\begin{ideabox}
\textbf{Зачем прошивать из командной строки?} Можно
использовать STM32CubeIDE~--- <<большую>> IDE от ST.
Она \emph{удобна}, но \emph{скрывает} детали: где
лежит бинарник, какие флаги у компилятора, что делает
линкер. Если вы хотите \emph{понимать}, что происходит~---
командная строка даёт \emph{полный} контроль.
\end{ideabox}

\subsection{Что вы узнаете}

\begin{itemize}
    \item Какие инструменты нужны и как их установить.
    \item Что такое ST-Link и SWD.
    \item Как убедиться, что ST-Link видит плату.
    \item Что делать, если в Flash лежит мусор.
    \item Как стереть Flash.
    \item Как записать и прочитать байты.
    \item Как подключиться отладчиком и посмотреть регистры.
    \item Как прошить первую программу~--- мигание LED.
\end{itemize}

\section{Инструменты}
\label{sec:flashing-tools}

Для работы с STM32 из командной строки нужно
\emph{пять} инструментов.

\subsection{ARM GCC: компилятор}

\begin{center}
\texttt{arm-none-eabi-gcc}
\end{center}

Это \textbf{кросс-компилятор}: работает на \emph{x86\_64} (ваш
компьютер), а генерирует код для \emph{ARM Cortex-M}
(STM32). <<\texttt{none-eabi}>> означает <<без операционной
системы, с ABI EABI>>~--- то есть для \emph{микроконтроллеров},
а не для полноценного Linux.

\textbf{Установка}:

\begin{minted}[frame=none, bgcolor=gray!5]{bash}
sudo apt install gcc-arm-none-eabi binutils-arm-none-eabi
\end{minted}

\textbf{Проверка}:

\begin{minted}[frame=none, bgcolor=gray!5]{bash}
arm-none-eabi-gcc --version
\end{minted}

\begin{keybox}
\textbf{Почему не обычный \texttt{gcc}?} Обычный \texttt{gcc}
компилирует код для \emph{вашего} компьютера~--- например,
для x86\_64. Код для ARM он не умеет~--- нужен
\emph{кросс-компилятор}. Это \emph{другой} GCC, собранный
под ARM.
\end{keybox}

\subsection{GDB: отладчик}

\begin{center}
\texttt{gdb-multiarch}
\end{center}

Это \textbf{отладчик}. Он умеет \emph{подключаться} к
STM32 через ST-Link, читать регистры, ставить точки
останова, выполнять код \emph{по шагам}.

\textbf{Установка}:

\begin{minted}[frame=none, bgcolor=gray!5]{bash}
sudo apt install gdb-multiarch
\end{minted}

\begin{warnbox}
\textbf{Грабля: пакет \texttt{gdb-arm-none-eabi} удалён.}
Начиная с \textbf{Debian 10}, пакет \texttt{gdb-arm-none-eabi}
\emph{отсутствует}. Если вы попробуете его установить~---
получите ошибку:
\begin{minted}{text}
Ошибка: Невозможно найти пакет gdb-arm-none-eabi
\end{minted}
\textbf{Решение:} установить \texttt{gdb-multiarch} и
создать симлинк:
\begin{minted}[frame=none, bgcolor=gray!5]{bash}
sudo ln -s /usr/bin/gdb-multiarch /usr/bin/arm-none-eabi-gdb
\end{minted}
После этого команда \texttt{arm-none-eabi-gdb} будет работать.
\end{warnbox}

\subsection{stlink-tools: работа с ST-Link}

\begin{center}
\texttt{st-info}, \texttt{st-flash}, \texttt{st-util}
\end{center}

Это \textbf{набор утилит} для работы с ST-Link~---
программатором STM32.

\textbf{Установка}:

\begin{minted}[frame=none, bgcolor=gray!5]{bash}
sudo apt install stlink-tools
\end{minted}

\textbf{Три утилиты}:

\begin{itemize}
    \item \texttt{st-info --probe}~--- \emph{обнаружить}
          ST-Link и показать информацию о чипе.
    \item \texttt{st-flash}~--- \emph{прошить} Flash,
          \emph{стереть}, \emph{прочитать}.
    \item \texttt{st-util}~--- \emph{GDB-сервер}~---
          подключается к ST-Link и транслирует запросы GDB.
\end{itemize}

\subsection{OpenOCD: альтернатива st-util}

\begin{center}
\texttt{openocd}
\end{center}

\textbf{OpenOCD} (Open On-Chip Debugger)~--- ещё один
GDB-сервер. Мы используем его, потому что \texttt{st-util 1.8.0}
\emph{несовместим} с \texttt{GDB 16.3}~--- читает Flash
\emph{неправильно}.

\textbf{Установка}:

\begin{minted}[frame=none, bgcolor=gray!5]{bash}
sudo apt install openocd
\end{minted}

\begin{notebox}
\textbf{Зачем два GDB-сервера?} \texttt{st-util}~---
\emph{родной} для stlink-tools, но старый. \texttt{OpenOCD}~---
более \emph{современный}, поддерживает много разных
адаптеров и целей. Если \texttt{st-util} работает~---
можно им пользоваться. Но для \emph{наших} задач
OpenOCD надёжнее.
\end{notebox}

\subsection{Сборка и git}

\begin{minted}[frame=none, bgcolor=gray!5]{bash}
sudo apt install build-essential pkg-config git
\end{minted}

\begin{itemize}
    \item \texttt{build-essential}~--- \texttt{gcc}, \texttt{make},
          \texttt{binutils} для \emph{хоста}.
    \item \texttt{pkg-config}~--- для поиска библиотек.
    \item \texttt{git}~--- для клонирования репозиториев.
\end{itemize}

\subsection{Итоговая команда установки}

\begin{minted}[frame=none, bgcolor=gray!5]{bash}
sudo apt install gcc-arm-none-eabi binutils-arm-none-eabi \
                 gdb-multiarch stlink-tools openocd \
                 build-essential pkg-config git
\end{minted}

\begin{keybox}
\textbf{Что мы \emph{не} устанавливали?} STM32CubeIDE,
STM32CubeMX, HAL-библиотеки. Для первой программы они
\emph{не нужны}. Мы напишем всё \emph{сами}~---
linker-скрипт, startup-код, мигание LED. Понимание
того, \emph{как} это работает~--- важнее, чем удобство.
\end{keybox}

\section{ST-Link: программатор}
\label{sec:flashing-stlink}

\textbf{ST-Link}~--- это \emph{программатор} (или
<<отладчик>>). Он подключается к компьютеру через USB
и к STM32 через \textbf{SWD}. Через него мы \emph{заливаем}
прошивку в Flash и \emph{отлаживаем} код.

\subsection{SWD: четыре провода}

\textbf{SWD} (Serial Wire Debug)~--- это \emph{протокол}
отладки ARM. Он использует \emph{всего четыре} провода:

\begin{itemize}
    \item \textbf{GND}~--- <<земля>>.
    \item \textbf{SWCLK}~--- тактовая линия.
    \item \textbf{SWDIO}~--- линия данных.
    \item \textbf{3.3V}~--- питание.
\end{itemize}

\begin{keybox}
\textbf{SWD~--- это <<два провода для отладки>>.}
SWCLK и SWDIO~--- тактовая и данные. По ним ST-Link
передаёт команды <<прочитай память>>, <<запиши байт>>,
<<останови процессор>>.
\end{keybox}

\subsection{Встроенный ST-Link на Nucleo}

\textbf{NUCLEO-F401RE} имеет \emph{встроенный} ST-Link~---
не нужен внешний программатор. Просто подключаете
плату по USB~--- и всё работает.

\begin{ideabox}
\textbf{Nucleo~--- это два устройства на одной плате.}
Одна половина~--- ST-Link (программатор). Другая~---
STM32F401 (целевой чип). Они соединены \emph{внутри}
платы по SWD. Наружу выходит только USB~---
и вы можете с ним работать, как с обычным ST-Link.
\end{ideabox}

\textbf{Проверка}:

\begin{minted}[frame=none, bgcolor=gray!5]{bash}
lsusb | grep -i st
\end{minted}

\textbf{Вывод} (наш случай):

\begin{minted}[frame=none, bgcolor=gray!5]{text}
Bus 001 Device 005: ID 0483:374b STMicroelectronics ST-LINK/V2.1
\end{minted}

\begin{itemize}
    \item \texttt{0483}~--- \emph{Vendor ID} STMicroelectronics.
    \item \texttt{374b}~--- \emph{Product ID} ST-Link/V2.1.
\end{itemize}

\subsection{Внешний ST-Link V2 для Blue Pill}

Blue Pill~--- это \emph{просто} микроконтроллер.
Программатора на нём \emph{нет}. Нужен внешний ST-Link V2.

\textbf{Подключение}:

\begin{center}
\begin{tabular}{ll}
ST-Link & Blue Pill \\
\hline
GND     & GND \\
SWCLK   & PA14 \\
SWDIO   & PA13 \\
3.3V    & 3.3V \\
\end{tabular}
\end{center}

\begin{warnbox}
\textbf{Не подавайте 5 В на SWD.} Это \emph{сожжёт} STM32.
Питание~--- только 3.3 В. Если сомневаетесь~--- измерьте
мультиметром.
\end{warnbox}

\section{Первое подключение: \texttt{st-info --probe}}
\label{sec:flashing-first-connect}

\textbf{Прежде чем} прошивать~--- убедимся, что ST-Link
\emph{видит} STM32. Это \emph{первый} шаг.

\begin{minted}[frame=none, bgcolor=gray!5]{bash}
st-info --probe
\end{minted}

\textbf{Вывод} (наш реальный случай):

\begin{minted}[frame=none, bgcolor=gray!5]{text}
Found 1 stlink programmers
  version:    V2J19S3
  serial:     066CFF484849785087031235
  flash:      524288 (pagesize: 16384)
  sram:       98304
  chipid:     0x433
  dev-type:   STM32F401xD_xE
\end{minted}

Разберём \emph{каждое} поле.

\subsection{Разбор вывода}

\begin{itemize}
    \item \texttt{Found 1 stlink programmers}~--- найдена
          \emph{одна} ST-Link. \textbf{Важно:} если здесь
          \texttt{2} или больше~--- у вас \emph{несколько}
          ST-Link подключены. Отключите лишние.
    \item \texttt{version: V2J19S3}~--- версия прошивки
          ST-Link. V2~--- ST-Link V2. J19~--- версия
          прошивки. S3~--- аппаратная ревизия.
    \item \texttt{serial: 066CFF...}~--- \emph{серийный
          номер} ST-Link. \textbf{Уникален.} Понадобится,
          если у вас \emph{несколько} ST-Link.
    \item \texttt{flash: 524288}~--- 512 КБ Flash
          (\SI{524288}{\byte} = 512 КБ).
    \item \texttt{pagesize: 16384}~--- размер страницы
          Flash~--- 16 КБ. \emph{Важно}: при стирании
          стирается \emph{вся} страница, а не отдельные
          байты.
    \item \texttt{sram: 98304}~--- 96 КБ SRAM.
    \item \texttt{chipid: 0x433}~--- идентификатор чипа.
          \texttt{0x433}~--- STM32F401.
    \item \texttt{dev-type: STM32F401xD\_xE}~--- точное
          название чипа.
\end{itemize}

\begin{keybox}
\textbf{Если \texttt{st-info --probe} показывает
\emph{не тот} чип}~--- значит, вы подключились
\emph{не к той} плате. Например, у вас есть Nucleo
(F401) и Blue Pill (F103), и вы хотите работать
с Nucleo, но подключились к Blue Pill.
\end{keybox}

\subsection{Грабля: две ST-Link одновременно}

\begin{warnbox}
\textbf{Мы наступили на эту граблю.} У нас была
\emph{вторая} ST-Link (для Blue Pill) подключена
\emph{одновременно} с Nucleo. Команды \texttt{st-info},
\texttt{st-flash}, \texttt{openocd} \emph{видели} только
\emph{одну} из них~--- не всегда \emph{ту}, которую
мы хотели. В результате мы получали \emph{странные}
данные из Flash.

\textbf{Решение:} отключите \emph{все} ST-Link, кроме
той, с которой работаете. Если вам \emph{нужны} обе~---
используйте флаг \texttt{-s <serial>} в \texttt{st-util}
или \texttt{st-flash}, чтобы указать \emph{конкретную}.
\end{warnbox}

\section{Nucleo и Mass Storage}
\label{sec:flashing-massstorage}

\textbf{Nucleo-F401RE}~--- это не только ST-Link и STM32.
Это ещё и \emph{виртуальная флешка}, которую видит
компьютер.

\subsection{Что такое Mass Storage}

Подключите Nucleo к компьютеру и откройте файловый
менеджер~--- вы увидите \emph{новый} диск под названием
\texttt{NUCLEO}. На нём будут \emph{файлы}:

\begin{itemize}
    \item \texttt{mbd.htm}~--- страница с информацией
          о плате.
    \item \texttt{System Volume Information}~--- служебный
          каталог Windows.
    \item \texttt{fseventsd}~--- служебный каталог macOS.
    \item \texttt{.Trashes}~--- корзина macOS.
    \item \texttt{.metadata\_never\_index}~--- маркер для
          Spotlight (macOS).
\end{itemize}

\begin{ideabox}
\textbf{Зачем это нужно?} \emph{Виртуальная флешка}~---
это <<drag-and-drop>> прошивка. Вы \emph{копируете}
файл \texttt{.bin} на диск \texttt{NUCLEO}~--- и
STM32 \emph{перепрошивается}. Никаких \emph{команд},
никаких \emph{инструментов}. Но мы будем использовать
\emph{командную строку}~--- для \emph{понимания}.
\end{ideabox}

\subsection{Грабля: \texttt{[Trash Info]} в Flash}

\begin{warnbox}
\textbf{Наш реальный случай.} Когда мы впервые
прочитали Flash через \texttt{st-flash read}, мы
увидели \emph{не} код ARM, а \emph{текст}:
\texttt{[Trash Info]}. Дальше~--- данные \emph{из} \emph{корзины}
macOS: \texttt{Path=System\%20Volume\%20Information},
\texttt{DeletionDate=2015-01-09T20:35:41}.

\emph{Что это было?} Заводская прошивка Nucleo
\emph{была повреждена} \emph{до} того, как плата
попала к нам. В неё \emph{случайно} попал файл из
корзины. При старте STM32 читал \emph{мусор} как
\emph{векторную таблицу}~--- и падал в \textbf{HardFault}.

\textbf{Решение:} \texttt{st-flash erase}~--- стереть
Flash полностью. После этого~--- прошить заново.
\end{warnbox}

\subsection{Как отключить Mass Storage (если мешает)}

Иногда \emph{виртуальная флешка} мешает работе~---
\textbf{GDB} \emph{путается} между интерфейсами.
Отключить её можно так:

\begin{enumerate}
    \item Отключите Nucleo от USB.
    \item Зажмите кнопку \textbf{USER} (синяя).
    \item Подключите USB (не отпуская USER).
    \item Подождите 2 секунды.
    \item Отпустите USER.
\end{enumerate}

\begin{notebox}
\textbf{Это} \emph{не всегда нужно}. Но если вы
видите \emph{странное} поведение~--- попробуйте.
\end{notebox}

\section{OpenOCD: альтернатива st-util}
\label{sec:flashing-openocd}

\textbf{OpenOCD}~--- это GDB-сервер. Он \emph{подключается}
к ST-Link и \emph{транслирует} команды GDB в SWD-запросы.

\subsection{Запуск}

\begin{minted}[frame=none, bgcolor=gray!5]{bash}
openocd -f interface/stlink.cfg -f target/stm32f4x.cfg
\end{minted}

\textbf{Вывод} (наш реальный случай):

\begin{minted}[frame=none, bgcolor=gray!5]{text}
Open On-Chip Debugger 0.12.0
Info : auto-selecting first available session transport "hla_swd".
Info : Listening on port 6666 for tcl connections
Info : Listening on port 4444 for telnet connections
Info : clock speed 2000 kHz
Info : STLINK V2J19M3 (API v2) VID:PID 0483:374B
Info : Target voltage: 3.236292
Info : [stm32f4x.cpu] Cortex-M4 r0p1 processor detected
Info : [stm32f4x.cpu] target has 6 breakpoints, 4 watchpoints
Info : starting gdb server for stm32f4x.cpu on 3333
Info : Listening on port 3333 for gdb connections
\end{minted}

\textbf{Что это значит}:

\begin{itemize}
    \item \texttt{Listening on port 3333}~--- сервер
          GDB работает на порту \textbf{3333}.
    \item \texttt{STLINK V2J19M3}~--- ST-Link \emph{найден},
          прошивка V2.J19.M3.
    \item \texttt{Target voltage: 3.236292}~--- напряжение
          на целевом STM32~--- \SI{3.24}{\volt}. Норма.
    \item \texttt{Cortex-M4 r0p1 processor detected}~---
          STM32F401 \emph{обнаружен}.
    \item \texttt{6 breakpoints, 4 watchpoints}~---
          \emph{аппаратные} точки останова.
\end{itemize}

\subsection{HardFault: что это значит}

\textbf{После} подключения GDB к OpenOCD, при попытке
\emph{остановить} процессор, мы видим:

\begin{minted}[frame=none, bgcolor=gray!5]{text}
[stm32f4x.cpu] halted due to debug-request,
current mode: Handler HardFault
xPSR: 0x01000003 pc: 0x7461500a msp: 0x61725438
\end{minted}

\textbf{HardFault}~--- это \emph{аппаратное} исключение.
Процессор \emph{встретил} что-то, что \emph{не может}
обработать: \emph{невалидную} инструкцию, \emph{обращение}
к \emph{несуществующей} памяти, \emph{деление} на ноль.

\textbf{В нашем случае} причина~--- Flash содержала
\emph{мусор}. Процессор \emph{читал} его как
\emph{инструкции} и \emph{падал} в HardFault.

\begin{keybox}
\textbf{HardFault~--- это симптом, а не причина.}
Если видите HardFault~--- ищите, \emph{почему} он
возник. В нашем случае~--- потому что Flash была
\emph{неправильной}. Решение~--- \texttt{st-flash erase}.
\end{keybox}

\subsection{PC и SP: почему мусор}

\begin{itemize}
    \item \texttt{pc: 0x7461500a}~--- это \emph{ASCII-код}:
          \texttt{0a 50 61 74} = \texttt{\textbackslash nPat}.
          То есть PC указывает на \emph{текст} в Flash.
    \item \texttt{msp: 0x61725438}~--- тоже \emph{ASCII}:
          \texttt{38 54 72 61} = \texttt{8Tra}.
    \item \texttt{xPSR: 0x01000003}~--- статус. Бит 3
          (\texttt{0x8})~--- \emph{прерывание}. То есть
          процессор \emph{сейчас} в обработчике HardFault.
\end{itemize}

\begin{ideabox}
\textbf{Как понять, что PC~--- это <<мусор>>?} Нормальный
PC для STM32F401~--- в диапазоне \texttt{0x08000000}–\texttt{0x0807FFFF}
(Flash) или \texttt{0x20000000}–\texttt{0x20017FFF} (SRAM).
Если PC \emph{вне} этих диапазонов~--- что-то не так.
\end{ideabox}

\section{Стирание Flash: \texttt{st-flash erase}}
\label{sec:flashing-erase}

\textbf{Прежде чем} писать что-то \emph{своё} в Flash~---
нужно \emph{стереть} то, что там \emph{лежит}. Даже если
кажется, что Flash <<пустая>>~--- там могут быть
\emph{остатки} старой прошивки.

\subsection{Зачем стирать}

\textbf{Три} причины:

\begin{enumerate}
    \item \textbf{Мусор.} Flash может содержать \emph{не то},
          что вы ожидаете~--- как в нашем случае
          с \texttt{[Trash Info]}.
    \item \textbf{Защита от случайного запуска.} После
          стирания Flash \emph{точно} пустая~---
          процессор \emph{не} выполнит случайный код.
    \item \textbf{Чистое состояние.} Мы \emph{знаем},
          что там \emph{все} \texttt{0xFF}~--- и
          можем \emph{проверить} результат записи.
\end{enumerate}

\subsection{Команда}

\begin{minted}[frame=none, bgcolor=gray!5]{bash}
st-flash erase
\end{minted}

\textbf{Вывод}:

\begin{minted}[frame=none, bgcolor=gray!5]{text}
st-flash 1.8.0
2026-10-15T07:50:38 INFO common.c: STM32F401xD_xE: 96 KiB SRAM, 512 KiB flash in at least 16 KiB pages.
Mass erasing..........
Mass erase completed successfully.
\end{minted}

\begin{keybox}
\textbf{Mass erase completed successfully}~--- это
<<всё стёрто>>. Теперь Flash содержит \emph{только}
байты \texttt{0xFF}.
\end{keybox}

\subsection{Проверка}

\begin{minted}[frame=none, bgcolor=gray!5]{bash}
st-flash read /tmp/nucleo-erased.bin 0x08000000 256
hexdump -C /tmp/nucleo-erased.bin | head -20
\end{minted}

\textbf{Вывод}:

\begin{minted}[frame=none, bgcolor=gray!5]{text}
00000000  ff ff ff ff ff ff ff ff  ff ff ff ff ff ff ff ff  |................|
*
00000100
\end{minted}

\begin{ideabox}
\textbf{Почему \texttt{0xFF}, а не \texttt{0x00}?} Потому
что \emph{стёртая} Flash NOR-типа \emph{всегда} содержит
\texttt{0xFF}. Запись~--- это <<перевод единиц в нули>>.
Стереть~--- значит <<вернуть все единицы>>. Это
\emph{аппаратная} особенность.
\end{ideabox}

\section{Запись и чтение: \texttt{st-flash write/read}}
\label{sec:flashing-write}

\textbf{Прежде чем} писать \emph{настоящую} программу~---
проверим, что \emph{запись/чтение} вообще \emph{работают}.

\subsection{Создаём тестовый файл}

Запишем \emph{восемь} байт:

\begin{minted}[frame=none, bgcolor=gray!5]{bash}
echo -e '\x00\x20\x00\x20\x00\x00\x00\x00' > /tmp/test.bin
\end{minted}

\begin{itemize}
    \item \texttt{00 20 00 20}~--- это \texttt{0x20002000}
          (в little-endian)~--- \emph{валидный} адрес в SRAM.
    \item \texttt{00 00 00 00}~--- это \texttt{0x00000000}.
\end{itemize}

\textbf{Почему именно так?} Потому что \emph{первые}
8 байт Flash~--- это \emph{векторная таблица}:
первые 4 байта~--- начальный \textbf{SP}, следующие
4~--- адрес \textbf{Reset Handler}. Мы <<эмулируем>>
векторную таблицу.

\subsection{Запись}

\begin{minted}[frame=none, bgcolor=gray!5]{bash}
st-flash write /tmp/test.bin 0x08000000
\end{minted}

\textbf{Вывод}:

\begin{minted}[frame=none, bgcolor=gray!5]{text}
st-flash 1.8.0
2026-10-15T23:05:47 INFO common.c: STM32F401xD_xE: 96 KiB SRAM, 512 KiB flash in at least 16 KiB pages.
file /tmp/test.bin md5 checksum: 64dafa57f4feea2243946d37802a7230, stlink checksum: 0x0000004a
2026-10-15T23:05:47 INFO common_flash.c: Attempting to write 9 (0x9) bytes to stm32 address: 134217728 (0x8000000)
2026-10-15T23:05:47 WARN common_flash.c: unaligned len 0x9 -- padding with zero
EraseFlash - Sector:0x0 Size:0x4000 -> Flash page at 0x8000000 erased (size: 0x4000)
2026-10-15T23:05:47 INFO flash_loader.c: Starting Flash write for F2/F4/F7/L4
2026-10-15T23:05:47 INFO flash_loader.c: Successfully loaded flash loader in sram
2026-10-15T23:05:47 INFO flash_loader.c: Clear DFSR
2026-10-15T23:05:47 INFO flash_loader.c: Clear CFSR
2026-10-15T23:05:47 INFO flash_loader.c: Clear HFSR
2026-10-15T23:05:47 INFO flash_loader.c: enabling 32-bit flash writes
2026-10-15T23:05:47 INFO common_flash.c: Starting verification of write complete
2026-10-15T23:05:47 INFO common_flash.c: Flash written and verified! jolly good!
\end{minted}

\subsection{Разбор вывода}

\begin{itemize}
    \item \texttt{md5 checksum: ...}~--- контрольная сумма
          файла. Если \emph{не} совпадёт с \texttt{stlink
          checksum}~--- \textbf{ошибка}.
    \item \texttt{Attempting to write 9 (0x9) bytes}~---
          \texttt{st-flash} прочитал \textbf{9} байт из
          файла (а не 8!). Почему? Потому что \texttt{echo}
          добавил \emph{новую строку} в конец.
    \item \texttt{unaligned len 0x9 -- padding with zero}~---
          длина \textbf{9} \emph{не кратна} 4. Flash
          пишется \emph{словами} по 32 бита~---
          \texttt{st-flash} \emph{дополнил} нулями до 12.
    \item \texttt{EraseFlash - Sector:0x0}~--- перед
          записью \texttt{st-flash} \emph{стёр} сектор 0
          (\SI{16}{\kilo\byte}).
    \item \texttt{Successfully loaded flash loader in sram}~---
          \texttt{st-flash} загрузил \emph{маленький
          загрузчик} в SRAM STM32. Этот загрузчик умеет
          \emph{писать} в Flash~--- само \emph{ядро} Cortex-M
          так не умеет.
    \item \texttt{Clear DFSR/CFSR/HFSR}~--- сброс регистров
          ошибок (Fault Status Registers). Профилактика.
    \item \texttt{enabling 32-bit flash writes}~--- включили
          запись \emph{словами}, а не байтами. \emph{Быстрее}.
    \item \texttt{Flash written and verified! jolly good!}~---
          \textbf{успех!} Данные \emph{записаны} и
          \emph{проверены}.
\end{itemize}

\begin{keybox}
\textbf{<<jolly good!>>~--- это фишка stlink-tools.}
Когда запись и верификация прошли успешно, утилита
пишет это \emph{по-английски}. Если увидели~---
всё \emph{точно} работает.
\end{keybox}

\subsection{Чтение обратно}

\begin{minted}[frame=none, bgcolor=gray!5]{bash}
st-flash read /tmp/test-read.bin 0x08000000 8
hexdump -C /tmp/test-read.bin
\end{minted}

\textbf{Вывод}:

\begin{minted}[frame=none, bgcolor=gray!5]{text}
00000000  00 20 00 20 00 00 00 00                           |. . ....|
00000008
\end{minted}

\textbf{Мы видим} \emph{ровно} те байты, которые записали.
\textbf{Запись/чтение работают}.

\begin{ideabox}
\textbf{Зачем это упражнение?} Оно \emph{доказывает},
что цепочка \texttt{компьютер → ST-Link → SWD → Flash}
работает. Если бы оно \emph{не} сработало~---
не было бы смысла писать \emph{настоящую} программу.
\end{ideabox}

\section{GDB: подключение и исследование}
\label{sec:flashing-gdb}

\textbf{GDB}~--- это отладчик. Через него мы \emph{видим},
что происходит \emph{внутри} STM32.

\subsection{Запуск OpenOCD}

\textbf{Терминал 1}:

\begin{minted}[frame=none, bgcolor=gray!5]{bash}
openocd -f interface/stlink.cfg -f target/stm32f4x.cfg
\end{minted}

\textbf{Ждём}:

\begin{minted}[frame=none, bgcolor=gray!5]{text}
Info : starting gdb server for stm32f4x.cpu on 3333
Info : Listening on port 3333 for gdb connections
\end{minted}

\begin{warnbox}
\textbf{Порт 3333, а не 4242.} \texttt{st-util} слушает
порт \textbf{4242}. \texttt{OpenOCD}~--- \textbf{3333}.
Не перепутайте!
\end{warnbox}

\subsection{Запуск GDB}

\textbf{Терминал 2}:

\begin{minted}[frame=none, bgcolor=gray!5]{bash}
arm-none-eabi-gdb
\end{minted}

\textbf{Внутри GDB}:

\begin{minted}[frame=none, bgcolor=gray!5]{text}
(gdb) target extended-remote 127.0.0.1:3333
\end{minted}

\textbf{Вывод}:

\begin{minted}[frame=none, bgcolor=gray!5]{text}
Remote debugging using 127.0.0.1:3333
warning: No executable has been specified and target does not support
determining executable automatically.  Try using the "file" command.
0x7461500a in ?? ()
\end{minted}

\begin{itemize}
    \item \texttt{Remote debugging using ...}~--- подключились.
    \item \texttt{warning: No executable ...}~--- это
          \emph{не} ошибка. GDB просто \emph{не знает},
          какая программа сейчас в STM32 (мы не указали
          файл). Он \emph{всё равно} может читать память
          и регистры.
    \item \texttt{0x7461500a in ?? ()}~--- PC сейчас
          \texttt{0x7461500a}. Это \emph{мусор}.
\end{itemize}

\subsection{\texttt{info registers}}

\begin{minted}[frame=none, bgcolor=gray!5]{text}
(gdb) info registers
\end{minted}

\textbf{Вывод} (наш случай):

\begin{minted}[frame=none, bgcolor=gray!5]{text}
r0             0x0                 0
r1             0x0                 0
r2             0x0                 0
r3             0x0                 0
r4             0x0                 0
r5             0x0                 0
r6             0x0                 0
r7             0x0                 0
r8             0x0                 0
r9             0x0                 0
r10            0x0                 0
r11            0x0                 0
r12            0x0                 0
sp             0x61725438          0x61725438
lr             0xfffffff9          -7
pc             0x7461500a          0x7461500a
xPSR           0x1000003           16777219
fpscr          0x0                 0
msp            0x61725438          0x61725438
psp            0x0                 0
primask        0x0                 0
basepri        0x0                 0
faultmask      0x0                 0
control        0x0                 0
\end{minted}

\subsection{Что здесь важно}

\begin{itemize}
    \item \texttt{r0}–\texttt{r12}~--- \emph{нули}.
          Это нормально~--- после сброса регистры обнуляются.
    \item \texttt{sp = 0x61725438}~--- <<стек-поинтер>>.
          Но \texttt{0x61725438}~--- это \emph{не} адрес
          в SRAM. Это \emph{мусор}.
    \item \texttt{lr = 0xfffffff9}~--- <<link register>>.
          Значение \texttt{-7}~--- тоже странно.
    \item \texttt{pc = 0x7461500a}~--- <<program counter>>.
          Мусор.
    \item \texttt{xPSR = 0x1000003}~--- статус. Бит 0
          (\texttt{0x1})~--- режим Thumb. Бит 1
          (\texttt{0x2})~--- бит \texttt{2} (<<precise
          fault>>).
    \item \texttt{msp = 0x61725438}~--- <<main stack pointer>>.
          Тоже мусор.
\end{itemize}

\begin{warnbox}
\textbf{Почему всё так странно?} Потому что STM32
\emph{сейчас} в HardFault~--- процессор пытается
обработать \emph{ошибку}, но не может. Регистры содержат
\emph{последние} значения, которые привели к ошибке.
\end{warnbox}

\subsection{Чтение Flash}

\begin{minted}[frame=none, bgcolor=gray!5]{text}
(gdb) x/8z 0x08000000
\end{minted}

\textbf{Команда \texttt{x}}~--- <<examine>> (исследовать).
Суффикс \texttt{/8z}~--- <<8 слов, шестнадцатерично>>.

\textbf{До стирания} мы видели:

\begin{minted}[frame=none, bgcolor=gray!5]{text}
0x8000000:	0x6172545b	0x49206873	0x5d6f666e	0x7461500a
0x8000010:	0x79533d68	0x6d657473	0x56303225	0x6d756c6f
\end{minted}

\textbf{Это}~--- ASCII: \texttt{[Tra}, \texttt{sh I},
\texttt{nfo]}, \texttt{\textbackslash nPat}, \texttt{h=Sy},
\texttt{stem}, \texttt{\%20V}, \texttt{olum}.

\begin{ideabox}
\textbf{Как <<увидеть>> текст в hex?} Возьмите байты
\emph{по отдельности} и посмотрите на ASCII-таблицу.
Например, \texttt{0x7461500a} в little-endian~---
это байты \texttt{0a 50 61 74}, то есть \texttt{\textbackslash n},
\texttt{P}, \texttt{a}, \texttt{t}. Вместе~---
\texttt{\textbackslash nPat}.
\end{ideabox}

\subsection{Чтение SRAM}

\begin{minted}[frame=none, bgcolor=gray!5]{text}
(gdb) x/4z 0x20000000
\end{minted}

\textbf{Вывод} (наш случай):

\begin{minted}[frame=none, bgcolor=gray!5]{text}
0x20000000:	0x8d89d84c	0xe29749ef	0x248280bc	0x755be5e7
\end{minted}

\textbf{После сброса} SRAM должна содержать \emph{нули}
или \emph{почти} нули. У нас~--- мусор. Это \emph{значит},
что процессор \emph{работал} до сбоя.

\section{Первая программа: мигание LED}
\label{sec:flashing-blink}

\textbf{Теперь} — самое интересное. Мы напишем
\emph{первую} программу — мигание LED на Nucleo.

\subsection{LED на Nucleo-F401RE}

На Nucleo-F401RE есть \emph{три} LED:

\begin{itemize}
    \item \textbf{LD1} (зелёный, у USB)~--- \emph{питание}.
          Управляется \emph{аппаратно}, не STM32.
    \item \textbf{LD2} (зелёный, у кнопки)~--- \emph{пользовательский}.
          Подключён к пину \textbf{PA5} STM32. \emph{Вот}
          его мы и будем мигать.
    \item \textbf{LD3} (красный, <<ST-Link>>)~--- \emph{ST-Link
          активен}. Управляется \emph{ST-Link}, не STM32.
\end{itemize}

\begin{keybox}
\textbf{LD2~--- единственный LED, управляемый STM32.}
Он подключён к PA5. Если подать на PA5 высокий уровень~---
LED горит. Если низкий~--- гаснет.
\end{keybox}

\begin{warnbox}
\textbf{Грабля: многие думают, что LD3~--- пользовательский.}
На Nucleo-F401RE \textbf{LD3}~--- это \emph{красный} LED
ST-Link. Он \emph{не} управляется STM32. \textbf{LD2}~---
\emph{зелёный}, подключён к PA5. \emph{Вот} его и мигаем.
\end{warnbox}

\subsection{Что нужно для прошивки}

Минимальная программа на C для STM32 состоит из \emph{трёх}
файлов:

\begin{itemize}
    \item \textbf{\texttt{linker.ld}}~--- <<карта памяти>>.
          Говорит линкеру, \emph{где} разместить код
          (\texttt{0x08000000}) и данные (\texttt{0x20000000}).
    \item \textbf{\texttt{startup.c}}~--- векторная таблица
          и обработчик сброса.
    \item \textbf{\texttt{blink.c}}~--- сам код мигания.
\end{itemize}

\subsection{linker.ld}

\begin{listing}[htbp]
\begin{minted}[frame=none, bgcolor=gray!5]{text}
ENTRY(Reset_Handler)

MEMORY
{
    FLASH (rx)  : ORIGIN = 0x08000000, LENGTH = 512K
    SRAM  (rwx) : ORIGIN = 0x20000000, LENGTH = 96K
}

/* Верхушка стека — конец SRAM */
_estack = ORIGIN(SRAM) + LENGTH(SRAM);

SECTIONS
{
    .text :
    {
        KEEP(*(.vector_table))
        *(.text*)
        *(.rodata*)
    } > FLASH

    /* Адрес, откуда копировать .data */
    _sidata = LOADADDR(.data);

    .data :
    {
        _sdata = .;
        *(.data*)
        _edata = .;
    } > SRAM AT > FLASH

    .bss :
    {
        _sbss = .;
        *(.bss*)
        *(COMMON)
        _ebss = .;
    } > SRAM
}
\end{minted}
\caption{\texttt{linker.ld}~--- карта памяти STM32F401}
\label{lst:flashing-linker}
\end{listing}

\textbf{Что здесь важного}:

\begin{itemize}
    \item \texttt{ORIGIN = 0x08000000}~--- Flash начинается
          \emph{здесь} (по документации STM32F4).
    \item \texttt{ORIGIN = 0x20000000}~--- SRAM начинается
          \emph{здесь}.
    \item \texttt{\_estack}~--- верхушка стека (\texttt{0x20018000}).
          Первое слово векторной таблицы.
    \item \texttt{.vector\_table}~--- секция для \emph{векторной
          таблицы}. Помечена \texttt{KEEP}, чтобы линкер
          \emph{не} выбросил её.
    \item \texttt{AT > FLASH}~--- \emph{данные} хранятся
          в Flash, но \emph{копируются} в SRAM при старте.
    \item \texttt{\_sidata}, \texttt{\_sdata}, \texttt{\_edata}~---
          символы для \texttt{startup.c}.
\end{itemize}

\subsection{startup.c}

\begin{listing}[htbp]
\begin{minted}[frame=none, bgcolor=gray!5]{c}
#include <stdint.h>

extern uint32_t _estack;
extern uint32_t _sidata, _sdata, _edata;
extern uint32_t _sbss, _ebss;
extern int main(void);

void Reset_Handler(void);
void Default_Handler(void);

__attribute__((section(".vector_table")))
const uint32_t vector_table[] = {
    (uint32_t)&_estack,          /* 0: Initial SP */
    (uint32_t)Reset_Handler,     /* 1: Reset */
    (uint32_t)Default_Handler,   /* 2: NMI */
    (uint32_t)Default_Handler,   /* 3: HardFault */
    (uint32_t)Default_Handler,   /* 4: MemManage */
    (uint32_t)Default_Handler,   /* 5: BusFault */
    (uint32_t)Default_Handler,   /* 6: UsageFault */
    0, 0, 0, 0,                  /* 7-10: Reserved */
    (uint32_t)Default_Handler,   /* 11: SVCall */
    (uint32_t)Default_Handler,   /* 12: Debug Monitor */
    0,                           /* 13: Reserved */
    (uint32_t)Default_Handler,   /* 14: PendSV */
    (uint32_t)Default_Handler,   /* 15: SysTick */
};

void Reset_Handler(void)
{
    /* Копируем .data из Flash в SRAM */
    uint32_t *src = &_sidata;
    uint32_t *dst = &_sdata;
    while (dst < &_edata) *dst++ = *src++;

    /* Обнуляем .bss */
    dst = &_sbss;
    while (dst < &_ebss) *dst++ = 0;

    main();

    while (1);
}

void Default_Handler(void)
{
    while (1);
}
\end{minted}
\caption{\texttt{startup.c}~--- векторная таблица и обработчик сброса}
\label{lst:flashing-startup}
\end{listing}

\textbf{Что здесь важного}:

\begin{itemize}
    \item \texttt{\#include <stdint.h>}~--- \textbf{обязательно}.
          Без него \texttt{uint32\_t} не определён~---
          компилятор выдаст \texttt{unknown type name 'uint32\_t'}.
    \item \texttt{vector\_table}~--- это \emph{первое},
          что читает STM32 после сброса. Первые 4 байта~---
          \emph{SP}, следующие 4~--- \emph{Reset\_Handler}.
    \item \texttt{Reset\_Handler}~--- копирует данные
          \emph{из} Flash \emph{в} SRAM (\texttt{.data}),
          обнуляет \texttt{.bss}, вызывает \texttt{main}.
    \item \texttt{Default\_Handler}~--- \emph{заглушка}
          для всех \emph{неиспользуемых} прерываний.
\end{itemize}

\begin{warnbox}
\textbf{Грабля: \texttt{unknown type name 'uint32\_t'}.}
Если забыть \texttt{\#include <stdint.h>}, компилятор
выдаст \emph{десятки} ошибок. \texttt{uint32\_t}~---
это \emph{не} встроенный тип, а \emph{тип} из
стандартной библиотеки C.
\end{warnbox}

\subsection{blink.c}

\begin{listing}[htbp]
\begin{minted}[frame=none, bgcolor=gray!5]{c}
#include <stdint.h>

/* Регистры STM32F401 */
#define RCC_AHB1ENR  (*(volatile uint32_t *)0x40023830)
#define GPIOA_MODER  (*(volatile uint32_t *)0x40020000)
#define GPIOA_ODR    (*(volatile uint32_t *)0x40020014)

static void delay(volatile uint32_t count)
{
    while (count--) __asm__ volatile ("nop");
}

int main(void)
{
    /* 1. Включаем тактирование GPIOA */
    RCC_AHB1ENR |= (1u << 0);

    /* 2. Настраиваем PA5 как выход */
    GPIOA_MODER &= ~(3u << 10);  /* сбрасываем биты 10-11 */
    GPIOA_MODER |=  (1u << 10);  /* 01 = output */

    /* 3. Мигаем */
    while (1) {
        GPIOA_ODR ^= (1u << 5);  /* инвертируем PA5 */
        delay(500000);
    }
}
\end{minted}
\caption{\texttt{blink.c}~--- мигание LED LD2 на Nucleo-F401RE}
\label{lst:flashing-blink}
\end{listing}

\textbf{Что здесь важного}:

\begin{itemize}
    \item Регистры~--- это \emph{адреса памяти}. Читая
          и записывая их, мы управляем периферией.
    \item \texttt{RCC\_AHB1ENR}~--- включает тактирование
          (<<часы>>) для порта GPIOA. \emph{Без этого}
          GPIOA не работает.
    \item \texttt{GPIOA\_MODER}~--- режим пина. Для PA5~---
          \emph{выход}.
    \item \texttt{GPIOA\_ODR}~--- состояние пина (0 или 1).
          \texttt{\^{}=}~--- \emph{инвертировать}.
\end{itemize}

\subsection{Makefile}

\begin{listing}[htbp]
\begin{minted}[frame=none, bgcolor=gray!5]{make}
CC      = arm-none-eabi-gcc
OBJCOPY = arm-none-eabi-objcopy
SIZE    = arm-none-eabi-size
CFLAGS  = -mcpu=cortex-m4 -mthumb -nostdlib -ffreestanding \
          -fno-builtin -O2 -g3 -Wall -Wextra

all: blink.bin

blink.elf: blink.c startup.c linker.ld
    $(CC) $(CFLAGS) -T linker.ld -o $@ blink.c startup.c
    $(SIZE) $@

blink.bin: blink.elf
    $(OBJCOPY) -O binary $< $@

flash: blink.bin
    st-flash write blink.bin 0x08000000

clean:
    rm -f blink.elf blink.bin

.PHONY: all flash clean
\end{minted}
\caption{\texttt{Makefile}~--- сборка и прошивка мигания LED}
\label{lst:flashing-makefile}
\end{listing}


\begin{warnbox}
\textbf{Табуляция в Makefile.} Команды \emph{обязаны}
начинаться с \textbf{табуляции} (\texttt{\^{}I}), а не
с \emph{пробелов}. Если вставить код из книги — табуляция
может испортиться. Проверьте командой:
\begin{minted}{bash}
cat -A Makefile | grep '^I'
\end{minted}
\end{warnbox}

\textbf{Сборка и прошивка}:

\begin{minted}[frame=none, bgcolor=gray!5]{bash}
make clean
make
make flash
\end{minted}

\textbf{Результат} (наш реальный случай):

\begin{minted}[frame=none, bgcolor=gray!5]{text}
   text    data     bss     dec     hex filename
    228       0       0     228      e4 blink.elf

st-flash write blink.bin 0x08000000
...
Flash written and verified! jolly good!
\end{minted}

\textbf{LED LD2} \emph{мигает}. Первая программа работает.

\begin{ideabox}
\textbf{228 байт~--- это вся программа.} Без HAL,
без библиотек, без IDE. Только \emph{наш} код
и \emph{наше} понимание железа. \textbf{512 КБ}
Flash у STM32F401~--- использовано \textbf{0.04\%}.
Остальное~--- для осциллографа.
\end{ideabox}

\section{Уроки главы}
\label{sec:flashing-lessons}

\begin{enumerate}
    \item \textbf{Пять инструментов:} ARM GCC, GDB,
          stlink-tools, OpenOCD, make.
    \item \textbf{\texttt{gdb-arm-none-eabi} удалён}
          из Debian. Используйте \texttt{gdb-multiarch}.
    \item \textbf{\texttt{st-info --probe}~--- первый шаг.}
    \item \textbf{Две ST-Link~--- грабля.} Отключайте лишние.
    \item \textbf{\texttt{[Trash Info]} в Flash~---
          реальность.} \texttt{st-flash erase} лечит.
    \item \textbf{HardFault~--- симптом.} Ищите причину.
    \item \textbf{Минимальная программа:} linker.ld +
          startup.c + blink.c + Makefile.
    \item \textbf{Регистры~--- это адреса памяти.} Читая
          и записывая их, управляем железом.
    \item \textbf{LD2~--- пользовательский LED.} Не LD3!
          LD3~--- это ST-Link.
    \item \textbf{\texttt{\#include <stdint.h>} в startup.c~---
          обязательно.} Иначе \texttt{uint32\_t} не определён.
\end{enumerate}

\section{Чек-лист грабель}
\label{sec:flashing-pitfalls}

\begin{table}[htbp]
\centering
\caption{Грабли прошивки STM32}
\label{tab:flashing-pitfalls}
\small
\begin{tabularx}{\textwidth}{@{}c L L@{}}
\toprule
\textbf{№} & \textbf{Грабля} & \textbf{Решение} \\
\midrule
1 & \texttt{gdb-arm-none-eabi} не найден & \texttt{gdb-multiarch} + симлинк \\
2 & \texttt{st-info} видит не тот чип & Отключить лишние ST-Link \\
3 & Flash содержит мусор & \texttt{st-flash erase} \\
4 & \texttt{st-util} читает мусор & Использовать OpenOCD \\
5 & HardFault после прошивки & Проверить векторную таблицу \\
6 & \texttt{st-flash} ругается на unaligned & Нормально, padding нулями \\
7 & LED не мигает & Проверить тактирование GPIOA \\
8 & Программа <<зависает>> & Нет \texttt{while(1)} в конце \\
9 & \texttt{unknown type name 'uint32\_t'} & \texttt{\#include <stdint.h>} \\
10 & Мигает не тот LED & LD2, не LD3. LD3~--- ST-Link. \\
11 & \texttt{make: *** missing separator} & Табуляция в Makefile \\
\bottomrule
\end{tabularx}
\end{table}

\section{Упражнения}
\label{sec:flashing-exercises}

\begin{exercisebox}
\textbf{Упражнение 1.} Запишите \emph{десять} байт
в Flash по адресу \texttt{0x08001000}. Прочитайте
обратно. Совпадают ли байты?
\end{exercisebox}

\begin{exercisebox}
\textbf{Упражнение 2.} Подключитесь GDB через OpenOCD.
Выполните \texttt{x/16z 0x08000000}. Сколько слов
занимает векторная таблица до \emph{нашего} Reset
Handler?
\end{exercisebox}

\needspace{5\baselineskip}
\begin{exercisebox}
\textbf{Упражнение 3.} Измените \texttt{delay()} в
\texttt{blink.c} так, чтобы LED мигал \emph{быстрее}.
Сколько времени занимает \emph{одна} итерация цикла?
\end{exercisebox}

\begin{exercisebox}
\textbf{Упражнение 4.} Добавьте в \texttt{vector\_table}
\emph{пустые} обработчики для всех прерываний STM32F401.
Сколько их всего?
\end{exercisebox}

\section{Проверьте сами}
\label{sec:flashing-check}

\begin{enumerate}
    \item Установите инструменты (см. раздел 2.2).
    \item Подключите Nucleo по USB.
    \item Проверьте: \texttt{st-info --probe}.
    \item Сотрите Flash: \texttt{st-flash erase}.
    \item Прочитайте: \texttt{st-flash read /tmp/x.bin 0x08000000 256}.
    \item Соберите программу: \texttt{make}.
    \item Прошейте: \texttt{make flash}.
    \item Убедитесь, что LED LD2 мигает.
\end{enumerate}

\begin{ideabox}
\textbf{Если LED мигает~--- вы прошли путь!} От \emph{нуля}
до \emph{первой} работающей программы на STM32. Дальше~---
интереснее: мы будем \emph{захватывать} сигналы, \emph{формировать}
их, \emph{рисовать} на экране.
\end{ideabox}
